\documentclass[14pt, a4paper]{extarticle}
\usepackage[
    a4paper,
    top=25mm,
    bottom=25mm,
    left=20mm,
    right=20mm
]{geometry}

\usepackage[T2A]{fontenc}
\usepackage[utf8]{inputenc}
\usepackage[russian]{babel}

\usepackage{amsmath, amssymb, amsthm}

\usepackage{setspace}
\onehalfspacing
\linespread{1.3}
\pagestyle{empty}
\begin{document}

\newgeometry{left=20mm, right=20mm, top=20mm, bottom=20mm}

\begin{center}
    {\large ГЛАВА 1} \\[0.5em]
    {\large \textbf{Введение}}
\end{center}
\vspace{1em}

\addcontentsline{toc}{section}{ГЛАВА 1. Введение}

Актуальность булевых функций и их исследований обусловлена развитием методов криптографического анализа
и ростом мощностей вычислительных машин, а также увеличением количества людей,имеющих к этим мощностям доступ.
В настоящий момент для взлома ряда криптографических систем не является обязательным использование суперкомпьютера,
входящего в рейтинг TOP500.

Для построения систем, например высокоскоростных поточных шифров,
аппаратно наиболее эффективным решением остается регистр сдвига с линейной обратной связью (РСЛОС, он же LFSR).
Однако сам по себе регистр имеет тривиальную алгебраическую структуру и может быть взломан быстрее,
чем за полный перебор,
при помощи алгоритма Берлекэмпа--Мэсси (Berlekamp--Massey) на основе перехваченной гаммы,
которая будет иметь очень малый объем (значение объема перехваченной гаммы подробно описано в Главе 3).

Для удаления линейных зависимостей
и обеспечения требуемого уровня защищенности в архитектуру генераторов псевдослучайных последовательностей внедряются системы нелинейного усложнения,
основой которых являются булевы функции.
Ключевая задача состоит в том,
что к современным функциям предъявляется большой спектр криптографических требований,
порой противоречащих друг другу (подробнее данный компромисс параметров описан в параграфе 3.3),
и даже 1 неудовлетворительный показатель из целого набора может забраковать функицю.

Поиск и верификация таких булевых функций в условиях сверхэкспоненциального роста вариантов представляет собой сложную комбинаторную задачу.
Если количество булевых функций от 5 переменных составляет $2^{32}$ (порядка 4~млрд),
то от 6 переменных --- уже $2^{64}$ (эквивалентно примерно 18,45 миллиарда миллиардов вариантов).
Решение этой задачи требует специализированного программного и алгоритмического обеспечения,
обладающего высокой производительностью.

Объектом исследования являются булевы функции,
применяемые в качестве узлов нелинейного усложнения в поточных криптосистемах.
Предметом исследования выступают алгоритмы многокритериального поиска,
анализа и фильтрации криптографических параметров булевых функций,
а также методы программной генерации структурированных математических задач.
Направление работы заключается в оптимизации подходов к параллельному многокритериальному анализу пространств булевых функций малой размерности
и их интеграции с алгоритмами автоматического синтеза пошаговых математических разборов сложных криптографических свойств.
Ценность результатов исследования состоит в создании программного обеспечения с графическим интерфейсом пользователя,
которое может быть эффективно использовано как в научной среде,
так и для проведения семинарских занятий по дисциплинам дискретной математики и криптографии.
Методологическую основу составляют положения теории конечных полей,
теории вероятностей, алгебры логики, а также принципы организации параллельных вычислений.

Целью выпускной квалификационной работы является разработка программного обеспечения,
предназначенного для многокритериального поиска криптографически стойких булевых функций.
Для достижения поставленной цели в рамках исследования решается комплекс задач,
включающий в себя изучение математического аппарата анализа булевых функций,
исследование современных криптографических атак на поточные шифры для формализации критериев фильтрации,
а также проектирование вычислительного ядра системы на базе математического инструментария среды SageMath
с разработкой дополнительных алгоритмов.
Программный блок также охватывает проектирование параллельного вычислительного модуля
под управлением асинхронного интерфейса пользователя
и создание модуля автоматизированного контроля знаний,
обеспечивающего динамическую верстку индивидуальных вариантов заданий и решений к ним в формате системы компьютерной верстки \LaTeX.

\newpage
\noindent{\large \textbf{1.1 \hspace{0.5em} Математическое описание и фундаментальные способы задания булевых функций}}
\vspace{1em}
\addcontentsline{toc}{section}{1.1 Математическое описание и фундаментальные способы задания булевых функций}

Изучение криптографических свойств узлов нелинейного усложнения
в современных поточных шифрах требует строгого математического описания пространств аргументов и значений дискретных отображений.
Перед анализом криптографических критериев необходимо формализовать сам объект исследования.

\vspace{2em}
\noindent{\textbf{1.1.1 \hspace{0.5em}Теоретико-множественное описание и проблема комбинаторного взрыва}}
\vspace{1em}
\addcontentsline{toc}{section}{1.1.1 Теоретико-множественное описание и проблема комбинаторного взрыва}


В рамках дискретного анализа булева функция от $n$ переменных
определяется как отображение вида $f: \mathbb{F}_2^n \to \mathbb{F}_2$,
где семейство элементов $\mathbb{F}_2 = \{0, 1\}$ представляет собой конечное поле из двух элементов.
Область определения функции образует дискретное векторное пространство $\mathbb{F}_2^n$ мощностью $2^n$.
Полное число различных булевых функций от фиксированного числа аргументов $n$ определяется выражением $2^{2^n}$.
Этот сверхэкспоненциальный рост создает вычислительные трудности для полного перебора уже при $n \ge 6$.

Поскольку отбор функций должен производится по нескольким критериям,
возникает необходимость в отказе от стратегий слепого перебора в пользу
методов направленного сужения пространства поиска за счет предварительного отсечения заведомо неперспективных подклассов.

\vspace{2em}
\noindent{\textbf{1.1.2 \hspace{0.5em}Способ задания булевых функций через полином Жегалкина}}
\vspace{1em}
\addcontentsline{toc}{section}{1.1.2 Способы задания булевых функций через полином Жегалкина}

Наиболее информативной формой представления является алгебраическая нормальная форма (АНФ),
или полином Жегалкина (к менее информативным были отнесены таблица истинности и СДНФ).
Она представляет собой сумму по модулю два различных конъюнктивных мономов,
где каждая переменная входит в произведение исключительно в первой степени.
Переход от вектора значений к полиномиальной форме эквивалентен выполнению линейного преобразования над полем $\mathbb{F}_2$.
В обобщенном случае АНФ имеет вид:
\begin{equation*}
f(x_1, x_2, \dots, x_n) = \bigoplus_{1 \le i < j \le n} x_i x_j
\end{equation*}

Степень старшего монома в полиноме определяет алгебраическую степень функции ($\deg f$).
Алгебраическая степень ограничивает стойкость криптосистемы к методам алгебраического анализа,
поскольку низкая степень позволяет криптоаналитику построить систему легко разрешимых уравнений (подробнее данный аспект рассматривается в Главе 3).

\vspace{2em}
\noindent{\textbf{1.1.3 \hspace{0.5em}Структурная классификация и подклассы булевых функций}}
\vspace{1em}
\addcontentsline{toc}{section}{1.1.3 Структурная классификация и подклассы булевых функций}

В зависимости от структуры полинома Жегалкина и распределения значений,
булевы функции классифицируются на несколько фундаментальных подклассов.
Ниже рассмотрены ключевые из них:

\begin{enumerate}
    \item \textbf{Линейные и аффинные функции ($L$).}
    Функция $f(x_1, x_2, \dots, x_n)$ называется аффинной, если её АНФ не содержит нелинейных мономов, то есть имеет вид:
    \begin{equation*}
    f(x_1, x_2, \dots, x_n) = c_0 \oplus c_1x_1 \oplus c_2x_2 \oplus \dots \oplus c_nx_n,
    \end{equation*}
    где $c_i \in \mathbb{F}_2$.
    Если константа $c_0 = 0$, функция называется линейной.
    Аффинные функции полностью уязвимы к методам линейного криптоанализа, так как их поведение на любых наборах данных тривиально предсказуемо.

    \item \textbf{Симметрические функции ($S$).}
    Функция является симметрической, если её значение зависит исключительно от веса Хемминга входного набора, но не от порядка следования переменных. 
    Формально, для любой перестановки $\alpha$ индексов аргументов должно выполняться равенство:
    \begin{equation*}
    f(x_1, x_2, \dots, x_n) = f(x_{\alpha(1)}, x_{\alpha(2)}, \dots, x_{\alpha(n)})
    \end{equation*}
    Использование симметрических функций позволяет оптимизировать аппаратную реализацию шифров, однако накладывает определенные ограничения на их лавинные свойства.

    \item \textbf{Самодвойственные функции ($S_0$).}
    Функция называется самодвойственной, если при инвертировании всех входных переменных значение самой функции также меняется на противоположное:
    \begin{equation*}
    f(\neg x_1, \neg x_2, \dots, \neg x_n) = \neg f(x_1, x_2, \dots, x_n)
    \end{equation*}
    В терминах таблицы истинности это означает, что вектор значений такой функции симметричен относительно своей середины с точностью до инверсии.
\end{enumerate}

\vspace{2em}
\noindent{\textbf{1.2 \hspace{0.5em}Числовые, метрические и геометрические свойства бинарных отображений}}
\vspace{1em}
\addcontentsline{toc}{section}{1.2 Числовые, метрические и геометрические свойства бинарных отображений}

Представление функции в виде таблицы или полинома не дает ответа на вопрос,
насколько эффективно она способна противостоять методам дешифрования.
Для этого вводится система количественных оценок — метрические и числовые характеристики бинарных отображений.

\vspace{2em}
\noindent{\textbf{1.2.1 \hspace{0.5em}Вес Хемминга, Метрига Хемминга и уравновешенность}}
\vspace{1em}
\addcontentsline{toc}{section}{1.2 Числовые, метрические и геометрические свойства бинарных отображений}

Вес Хемминга булевой функции $f(x_1, x_2, \dots, x_n)$, обозначаемый как $w(f)$,
определяется как количество единиц в векторе её значений:
\begin{equation*}
w(f) = |\{\alpha \in \mathbb{F}_2^n : f(\alpha) = 1\}|
\end{equation*}

Булева функция от $n$ переменных называется уравновешенной, если её вес равен ровно половине мощности области определения, то есть:
\begin{equation*}
w(f) = 2^{n-1}
\end{equation*}

Это означает, что при подаче на вход функции случайного двоичного набора, вероятность получения на выходе нуля или единицы одинакова и равна $P(f=0) = P(f=1) = 0.5$.
Уравновешенность — обязательное требование в поточных шифрах.
Её отсутствие создает статистическое смещение,
что позволяет злоумышленнику провести эффективные криптоатаки (например, линейный или корреляционный анализ).
Такие функции откидываются в первую очередь, как криптографически слабые
(впрочем, при надобности пользователь может эту проверку в программе отключить).

Для двух произвольных двоичных векторов $\alpha$ и $\beta$
из пространства $\mathbb{F}_2^n$ расстояние Хемминга определяется как число несовпадающих разрядов на соответствующих позициях,
выражаемое как вес их покомпонентной суммы
\begin{equation*}
d_H(\alpha, \beta) = w(\alpha \oplus \beta) = \sum_{i=1}^{n} (\alpha_i \oplus \beta_i)
\end{equation*}

К примеру, это применяется при вычислении нелинейности булевой функции (которая подробно исследуется в Главе 3).
Нелинейность определяется как минимальное расстояние Хемминга от функции до множества всех аффинных функций.
Высокая нелинейность усложняет линейную аппроксимацию шифра.

\vspace{2em}
\noindent{\textbf{1.2.2 \hspace{0.5em}Геометрия $n$-мерного гиперкуба и связь со SAC}}
\vspace{1em}
\addcontentsline{toc}{section}{1.2.2 Геометрия $n$-мерного гиперкуба и связь со SAC}

Сферой Хемминга радиуса $k$ с центром в точке $\alpha$ называется множество вершин, находящихся от неё на расстоянии ровно $k$:
\begin{equation*}
S_k(\alpha) = \{\beta \in \mathbb{F}_2^n : d_H(\alpha, \beta) = k\}
\end{equation*}

Мощность такой сферы равна $\binom{n}{k}$.
Шаром радиуса $k$ с веденной на нем метрикой Хемминга называется объединение всех сфер от нулевого до $k$-го радиуса,
включающее в себя все вершины, удаленные от центра не более чем на $k$ шагов:
\begin{equation*}
B_k(\alpha) = \{\beta \in \mathbb{F}_2^n : d_H(\alpha, \beta) \le k\} = \bigcup_{i=0}^{k} S_i(\alpha) \subset \mathbb{Z}_2^n
\end{equation*}

Если зафиксировать вершину $\alpha$ и рассмотреть её ближайшую окрестность — сферу $S_1(\alpha)$ 
состоящую из $n$ смежных вершин,
то инверсия любого одного бита входного вектора означает переход из центра сферы в одну из этих $n$ точек.
Характер изменения значений функции $f$ при таких переходах определяет её лавинные свойства
и соответствие строгому лавинному критерию (SAC — Strict Avalanche Criterion, детально рассматриваемому в пункте 3.4),
Если функция удовлетворяет SAC, то при переходе из $\alpha$ в любую точку сферы $S_1(\alpha)$
значение её выхода меняется с вероятностью $0.5$, что исключает проведение эффективных разностных криптоатак.

\vspace{2em}
\noindent{\textbf{1.3 \hspace{0.5em}Производная булевых функций по направлению}
\vspace{1em}
\addcontentsline{toc}{section}{1.3 Производная булевых функций по направлению}

Для бинарных отображений можно корректно ввести понятие дискретной производной по направлению.
Дискретная производная функции $f(x)$ по направлению вектора $u$ обозначается как $(D_u f)(x)$
и определяется как:
\begin{equation*}
(D_u f)(x) = f(x \oplus u) \oplus f(x), \quad x \in V_n
\end{equation*}

Математический смысл этого оператора аналогичен непрерывной производной:
функция $(D_u f)(x)$ принимает значение 1 в тех точках $x$,
где сдвиг аргумента на вектор $u$ приводит к инверсии значения самой функции, и 0 там,
где значение функции остается неизменным.
В дискретном пространстве характер распределения значений производной $(D_u f)(x)$
позволяет оценивать "хаотичность" и нелинейность отображения.
\end{document}